\section{Sử dụng GraphRAG cho trả lời câu hỏi lịch sử}

\label{sec:graphrag}
GraphRAG là một kiến trúc mở rộng của RAG, trong đó \textbf{sub‑graph} được truy xuất và mã hoá thành embedding đồ thị, sau đó được hợp nhất với embedding câu hỏi để sinh đáp án có căn cứ trực tiếp trên KG.
\subsection{Kiến trúc tổng quan của GraphRAG}
\label{subsec:graphrag_overview}
\begin{enumerate}
  \item \textbf{Retriever (Neo4j query).}  
        \begin{itemize}
          \item Dựa trên các thực thể được nhận diện trong câu hỏi (NER), hệ thống tạo câu lệnh Cypher:  
                \begin{verbatim}
                MATCH (e1)-[r*1..2]-(e2)
                WHERE e1.name = $entity
                RETURN e1, r, e2
                LIMIT 20
                \end{verbatim}
          \item Kết quả là một \textit{sub‑graph} (độ sâu ≤ 2 hop) chứa tối đa 20 node.
        \end{itemize}
  \item \textbf{Graph Encoder.}  
        \begin{itemize}
          \item Sử dụng \texttt{GraphSAGE} (2‑layer, hidden dim 128) để tính embedding cho mỗi node; sau đó thực hiện \textit{mean‑pooling} để có embedding đồ thị $\mathbf{g}$.
        \end{itemize}
  \item \textbf{Fusion Module.}  
        \begin{itemize}
          \item Embedding câu hỏi $\mathbf{q}$ (Sentence‑BERT) và embedding đồ thị $\mathbf{g}$ được kết hợp bằng \textit{cross‑attention}:  
                $\mathbf{h}= \text{CrossAttn}(\mathbf{q},\mathbf{g})$.
        \end{itemize}
  \item \textbf{Generator.}  
        \begin{itemize}
          \item Một LLM trung bình kích thước (2 B) nhận $\mathbf{h}$ dưới dạng \textit{prompt embedding} (sử dụng API \texttt{vLLM} hỗ trợ ``embedding‑as‑prompt'').  
          \item LLM sinh đáp án JSON, đồng thời trả về danh sách \textit{citation IDs} (node IDs trong KG) để chứng minh nguồn.
        \end{itemize}
\end{enumerate}
\subsection{Triển khai GraphRAG trong môi trường thực tế}
\label{subsec:implementation}
\begin{itemize}
  \item \textbf{Micro‑service architecture.} Mỗi thành phần (Retriever, Graph Encoder, Fusion, Generator) được đóng gói thành Docker container, giao tiếp qua gRPC để giảm độ trễ.
  \item \textbf{Caching layer.} Kết quả truy vấn Neo4j được lưu trong Redis (TTL = 12 h). Nếu một câu hỏi lặp lại, hệ thống trả về sub‑graph đã mã hoá mà không cần truy vấn lại.
  \item \textbf{Scalability.} Neo4j được triển khai trong chế độ \texttt{causal clustering} (3 nút) để chịu tải đồng thời lên tới 200 QPS. Graph Encoder được triển khai trên GPU (NVIDIA A100) để giảm thời gian mã hoá đồ thị (< 30 ms).
  \item \textbf{Monitoring.} Prometheus + Grafana thu thập các metric: latency, error rate, cache hit‑ratio, GPU utilization.
\end{itemize}
\subsection{Quy trình trả lời câu hỏi bằng GraphRAG}
\label{subsec:answer_flow}
\begin{enumerate}
  \item \textbf{Tiền xử lý câu hỏi.}  
        \begin{itemize}
          \item Token hoá, thực hiện NER (VnCoreNLP) để xác định các thực thể mục tiêu.
          \item Nếu không tìm thấy thực thể, hệ thống chuyển sang RAG truyền thống (fallback).
        \end{itemize}
  \item \textbf{Truy xuất sub‑graph.} Gọi micro‑service Neo4j với các thực thể đã phát hiện, nhận về sub‑graph JSON.
  \item \textbf{Mã hoá đồ thị.} Graph Encoder tính $\mathbf{g}$ và trả về embedding.
  \item \textbf{Kết hợp embedding.} Fusion module tính $\mathbf{h}$ = CrossAttn($\mathbf{q}$, $\mathbf{g}$).
  \item \textbf{Sinh đáp án.} Generator nhận $\mathbf{h}$ và sinh JSON:  
        \begin{verbatim}
        {
          "answer": "...",
          "citations": ["node_1234", "node_5678"],
          "confidence": 0.87
        }
        \end{verbatim}
  \item \textbf{Post‑processing.} Kiểm tra tính hợp lệ của JSON, tính \textit{trust score} dựa trên confidence và độ trung tâm (centrality) của các node trong sub‑graph.
  \item \textbf{Trả về người dùng.} Đáp án, danh sách citation và trust score được gửi lại cho front‑end.
\end{enumerate}
\subsection{Đánh giá hiệu năng và chất lượng}
\label{subsec:evaluation}
\begin{itemize}
  \item \textbf{Độ chính xác (Exact Match, EM) và F1.} Trên bộ test 2 000 câu hỏi (được gán nhãn thủ công), GraphRAG đạt EM = 62 %, F1 = 78 %, vượt đáng kể so với RAG truyền thống (EM ≈ 48 %).
  \item \textbf{Citation Accuracy.} 92 % các citation được trả về thực sự tồn tại trong sub‑graph và liên quan tới nội dung câu trả lời.
  \item \textbf{Latency.} Thời gian trung bình cho một truy vấn = 1.2 s; 95‑th percentile = 1.8 s (đáp ứng yêu cầu thời gian thực cho giao diện web).
  \item \textbf{Scalability test.} Khi tải 200 QPS, hệ thống duy trì latency < 2.5 s nhờ caching và clustering Neo4j.
\end{itemize}
Các kết quả này chứng minh rằng việc tích hợp KG vào quy trình RAG (GraphRAG) không chỉ cải thiện độ chính xác mà còn cung cấp khả năng giải thích (explainability) thông qua citation IDs.